\documentclass[10pt,a4paper]{amsart}
\usepackage[margin=19mm]{geometry}
\usepackage{amsmath,amssymb,mathtools}
\usepackage[scr=rsfs,cal=euler]{mathalfa}
\usepackage{xfrac}
\usepackage[unicode,hidelinks,bookmarks=false]{hyperref}
\newcommand{\ie}{i.e.\ }
\newcommand{\cf}{cf.\ }
\newcommand{\resp}{respectively\ }
\newcommand{\Subsection}{\subsection}
\providecommand{\nobreakdash}{\nobreak\hbox{-}}
\setlength{\emergencystretch}{2em}
\multlinegap0pt
\usepackage{amssymb}
\usepackage{mathtools}
\usepackage{xcolor}
\usepackage{tikz}
\usepackage[shortlabels]{enumitem}
\newtheorem{lemma}{Lemma}
\newcommand{\symb}[2]{\left(\frac{#1}{#2}\right)}
\title{Squares, Three Fleas, Sporadic Sets, and Squares}
\author{Giedrius Alkauskas}
\date{Source: arXiv:2508.05347v3 (CC BY 4.0)}
\begin{document}
\maketitle

\noindent\emph{Source and licence.} Squares, Three Fleas, Sporadic Sets, and Squares. Version arXiv:2508.05347v3 (CC BY 4.0), distributed by arXiv under the Creative Commons Attribution 4.0 International licence, http://creativecommons.org/licenses/by/4.0/, as recorded in the arXiv licence metadata for this exact version and shown on the abstract page https://arxiv.org/abs/2508.05347v3. Full licence notice: \texttt{licenses/arxiv-2508.05347-CC-BY-4.0.txt}.\par\smallskip

\noindent\textit{Editorial scope.} Counted unit: the article's lemma lem:ar (source0.tex lines 371--375). The article's complete original proof of it is delivered with it (source0.tex lines 377--410, one \texttt{proof} environment of 34 lines). No mathematics is written, repaired or re-derived: the statement and the proof are the article's own lines, in the article's own notation, and no retained line is substituted or re-flowed.  No further source-local context is needed: every cross-reference of every retained line resolves inside the two delivered spans.  Nothing was omitted from any retained span, so the package carries no omission record.  No retained line contains a citation, so the wrapper carries no bibliography and no bibliography entry.  No retained line contains a figure environment, an image-inclusion command or drawing code, so no image is delivered and the package declares no asset dependency.  Editorial formatting only, in addition: an AMS \texttt{amsart} preamble that restates the article's own declarations and macros used by the retained lines, namely the environments \texttt{lemma} on separate counters, together with the standard packages those lines need.  The retained environments print in this document as Lemma 1 in order of appearance, whereas the article prints the same items with the numbering of its own full document.  The complete native source of the article as distributed in the arXiv e-print for this exact version is bundled unchanged as \texttt{sources/182855/source0.tex}.
\begin{lemma}\label{lem:ar}
	Suppose that a flea state $(a,b,c;s)$ has primitive associated Descartes
	quadruple $Q$ and that $|s|=2w^2$, $w\ge1$. If exactly one of $a,b,c$ is odd,
	denote it by $n$. Then $\chi_Q(n)=1$.
\end{lemma}

\begin{proof}
	Relabel so that $n=a$, and put $R=a+b$, $T=a+c$. 
	The squared side lengths \(R=a+b\) and \(T=a+c\) are positive and, since \(b,c\) are even, odd. Now,
	\[
	RT=(a+b)(a+c)=a^2+s^2=n^2+4w^4.
	\]
Choose (as allowed by CRT) an integer $z$ such that	$z\equiv w\pmod{R^2}$, $z\equiv1\pmod p$ for every prime $p\mid n$ with $p\nmid R$. We shall manufacture an odd integer \(E\), coprime to \(n\), which reproduces the local character of \(n\) at every prime dividing \(n\), but whose own prime divisors force \(n\) to be a quadratic residue. Thus, define
	\[
	E=\frac{n^2+4z^4}{R}.
	\]
This is a positive odd integer: indeed,	$n^2+4z^4=RT+4(z^4-w^4)$, and $R\mid z^4-w^4$. We first identify the local character of $n$. Let $p\mid n$. If $p\nmid R$, then $E\equiv4R^{-1}\pmod p$, and therefore $\symb{E}{p}=\symb{R}{p}$. Since $R=n+b\equiv b\pmod p$, this is precisely the local character of $n$.
	
Suppose instead that $p\mid R$. From $RT=n^2+4w^4$ we obtain $p\mid w$. Moreover $p\nmid T$: otherwise $p$ would divide $a,b,c,s$, and hence also the fourth Descartes coordinate, contrary to primitivity. Since $z\equiv w\pmod{R^2}$,
	\[
	E-T=4\,\frac{z^4-w^4}{R}
	\]
is divisible by $R$, and hence $E\equiv T\pmod p$. Thus $\symb{E}{p}=\symb{T}{p}$, again the local character of $n$, since $T=n+c\equiv c\pmod p$. In either case $p\nmid E$. Hence $\gcd(E,n)=1$, and the definition of the character gives
	\begin{equation}
	\chi_Q(n)=\symb{E}{|n|}.                     
	\label{reco}
	\end{equation}
	It remains to evaluate this Jacobi symbol. Let $q$ be any prime divisor
	of $E$. Then $q\nmid n$, and $n^2+4z^4\equiv0\pmod q$. Also $q\nmid z$, for otherwise $q\mid n$. Consequently
	\[
	i=\frac{|n|}{2z^2}\pmod q
	\]
is defined and satisfies $i^2\equiv-1\pmod q$. Thus $q\equiv1\pmod4$. Moreover, since $(1+i)^2=2i$, $|n|\equiv2iz^2=\bigl((1+i)z\bigr)^2\pmod q$, and therefore $\symb{|n|}{q}=1$. Every prime divisor of the odd integer $E$ is therefore $1\pmod4$, so $E\equiv1\pmod4$. Quadratic reciprocity and \eqref{reco} now give
	\[
	\chi_Q(n)
	=\symb{E}{|n|}
	=\symb{|n|}{E}
	=1.
	\]
\end{proof}

\end{document}
